\chapter{Literature Review and Research Gaps}
\label{chap:literature_review}

Computational models driven by Machine Learning (ML) have emerged as essential assets for clinical data interpretation and diagnostic forecasting. Academic inquiries have consistently shown that supervised learning frameworks can achieve high precision in identifying conditions from structured medical records. However, a persistent gap exists between reported benchmarks and real-world clinical applicability.

\section{Machine Learning in Symptom-Based Diagnosis}
\label{sec:ml_in_healthcare}
Automated disease classification has evolved from rule-based expert systems (such as MYCIN and INTERNIST-I) toward statistical learning and deep neural networks. Multiple studies have applied Decision Trees, Random Forests, Support Vector Machines (SVM), and Gradient Boosted Trees (XGBoost) to tabular symptom-disease profiles. However, existing literature overwhelmingly reports near-perfect accuracies without investigating data redundancy or testing robustness against missing features.

\section{The Pitfall of Data Leakage in Healthcare AI}
\label{sec:data_leakage}
Data leakage occurs when information from outside the training dataset is inadvertently utilized to train the model~\cite{kapoor2023leakage}. In medical diagnostics, when synthetic benchmarks duplicate symptom vectors across train and validation splits, standard cross-validation grossly overestimates model generalization. Models learn to act as lookup tables for recurring synthetic vectors, leading to catastrophic failure when deployed on real patients presenting novel symptom combinations.

\section{Probability Calibration in Clinical Decision Support}
\label{sec:probability_calibration}
A clinical decision-support tool must express calibrated confidence~\cite{guo2017calibration}. If a classifier assigns a 90\% probability to Dengue, empirical clinical observations should confirm Dengue in approximately 90\% of such cases. Modern non-linear classifiers, particularly deep neural networks and boosted trees, often suffer from overconfidence. Post-hoc calibration via Platt scaling~\cite{platt1999probabilistic} and Isotonic regression~\cite{zadrozny2002transforming} is critical to transform raw scores into actionable clinical probabilities.

\section{Explainability in Medical Artificial Intelligence}
\label{sec:xai_literature}
Rudin~\cite{rudin2019stop} argues that for high-stakes decisions, interpretable models should be preferred over post-hoc explanations of black-box systems. In practice, the clinical deployment of ML-based diagnostic tools requires model transparency to foster physician trust. TreeSHAP~\cite{lundberg2017unified} provides exact and computationally efficient Shapley value explanations for tree-based models, enabling alignment between algorithmic reasoning and established medical knowledge.

\section{Research Gaps Addressed by This Work}
\label{sec:research_gaps}
Synthesizing the literature, the following critical gaps are identified and addressed in this thesis:
\begin{itemize}
    \item No prior published work has conducted a forensic audit of the Kaggle Disease Prediction dataset to quantify its synthetic duplication rate and train-test leakage severity.
    \item Robustness profiling under progressive symptom sparsity ($s \in [1, 10]$ active symptoms) with multi-candidate differential diagnostic evaluation remains absent from the literature.
    \item The interaction between feature selection method and leakage-aware evaluation under Protocol C has not been studied in the multi-disease symptom prediction setting.
\end{itemize}
